\section{Formula \& Key Results Reference}
%==================================================

\subsection{Ch1: Definite integrals, FTC, substitution, improper integrals}

\subsubsection{FTC / Leibniz / MVT}
\begin{keyformula}[title={Fundamental Theorem of Calculus and Leibniz Rules}]
\begin{align*}
&\textbf{FTC I:}\quad f\in C[a,b],\quad F(x):=\int_a^x f(t)\,dt
\quad\Rightarrow\quad F'(x)=f(x).\\
&\textbf{FTC II:}\quad f\in C[a,b],\quad F'=f
\quad\Rightarrow\quad \int_a^b f(x)\,dx=F(b)-F(a).\\
&\textbf{Variable limits:}\quad
\frac{d}{dx}\int_{g(x)}^{h(x)} f(t)\,dt
=f(h(x))h'(x)-f(g(x))g'(x).\\
&\textbf{Parameter-dependent integrand:}\quad
\frac{d}{dx}\int_{g(x)}^{h(x)} F(x,t)\,dt
=F(x,h(x))h'(x)-F(x,g(x))g'(x)+\int_{g(x)}^{h(x)}F_x(x,t)\,dt.
\end{align*}
\end{keyformula}

\begin{keyformula}[title={Mean Value Theorem for Integrals and Average Value}]
\[
f\in C[a,b]\Rightarrow \exists c\in(a,b):\quad
\int_a^b f(x)\,dx=f(c)(b-a),
\qquad
f_{\rm avg}=\frac{1}{b-a}\int_a^b f(x)\,dx.
\]
\end{keyformula}

\begin{examtip}
For a lower variable limit, the contribution is negative.  A quick check is
\[
\int_{g(x)}^{h(x)}f(t)\,dt=\int_a^{h(x)}f(t)\,dt-\int_a^{g(x)}f(t)\,dt.
\]
Always use a dummy variable inside the integral.
\end{examtip}

\subsubsection{$u$-substitution}
\begin{keyformula}[title={Substitution, Indefinite and Definite Forms}]
\begin{align*}
&u=g(x),\quad du=g'(x)\,dx
\quad\Rightarrow\quad
\int f(g(x))g'(x)\,dx=\int f(u)\,du.\\
&\int_a^b f(g(x))g'(x)\,dx=\int_{g(a)}^{g(b)} f(u)\,du.
\end{align*}
\end{keyformula}

\begin{commonmistake}
For a definite integral, either change the bounds and stay in the new variable,
or keep the old bounds and back-substitute.  Do not do both.
\end{commonmistake}

\subsubsection{Improper integrals}
\begin{keyformula}[title={Improper Integral Definitions}]
\begin{align*}
&\int_a^{\infty} f(x)\,dx:=\lim_{t\to\infty}\int_a^t f(x)\,dx,
\qquad
\int_{-\infty}^{b} f(x)\,dx:=\lim_{t\to-\infty}\int_t^b f(x)\,dx.\\
&\int_a^b f(x)\,dx:=\lim_{t\to a^+}\int_t^b f(x)\,dx
\quad\text{if }f\text{ is improper at }a,\quad
\text{similarly at }b.\\
&\text{If }f\text{ is improper at }c\in(a,b),\quad
\int_a^b f=\int_a^c f+\int_c^b f,
\quad\text{and both one-sided integrals must converge.}
\end{align*}
\end{keyformula}

\begin{keyformula}[title={Improper Integral Comparison and $p$-Tests}]
\begin{align*}
&\int_1^{\infty}x^{-p}\,dx\text{ converges}\Longleftrightarrow p>1,
\qquad
\int_0^1x^{-p}\,dx\text{ converges}\Longleftrightarrow p<1.\\
&0\le f\le g,
\quad \int g\text{ converges}\Rightarrow \int f\text{ converges};
\qquad
0\le g\le f,
\quad \int g\text{ diverges}\Rightarrow \int f\text{ diverges}.\\
&f,g>0,
\quad \lim_{x\to\alpha}\frac{f(x)}{g(x)}=c\in(0,\infty)
\Rightarrow \int f\text{ and }\int g\text{ have the same convergence behaviour.}
\end{align*}
\end{keyformula}

\begin{commonmistake}
Do not use the antiderivative value across an infinite endpoint or vertical
asymptote without first rewriting the integral as a limit.
\end{commonmistake}

\subsection{Ch2: Applications of integration}

\subsubsection{General setup principle}
\begin{examtip}
Before computing, choose the variable, mark the interval, split at intersections
or singularities, write the slice formula, and only then integrate.  Most errors
in application questions occur in the setup, not the antiderivative.
\end{examtip}

\subsubsection{Area}
\begin{keyformula}[title={Area Between Curves}]
\[
A=\int_a^b\bigl(y_{\rm top}(x)-y_{\rm bot}(x)\bigr)\,dx,
\qquad
A=\int_c^d\bigl(x_{\rm right}(y)-x_{\rm left}(y)\bigr)\,dy.
\]
\end{keyformula}

\begin{commonmistake}
A signed integral is not automatically geometric area.  Split at crossing points
or sign changes before taking area.
\end{commonmistake}

\subsubsection{Volumes}
\begin{keyformula}[title={Slicing, Washers, and Shells}]
\begin{align*}
&\textbf{Slicing:}\quad V=\int A(x)\,dx\quad\text{or}\quad V=\int A(y)\,dy.\\
&\textbf{Washers about }y=k:\quad
V=\pi\int_a^b\left(R_{\rm out}(x)^2-R_{\rm in}(x)^2\right)\,dx.\\
&\textbf{Washers about }x=k:\quad
V=\pi\int_c^d\left(R_{\rm out}(y)^2-R_{\rm in}(y)^2\right)\,dy.\\
&\textbf{Shells about }x=k:\quad
V=2\pi\int_a^b |x-k|\bigl(y_{\rm top}(x)-y_{\rm bot}(x)\bigr)\,dx.\\
&\textbf{Shells about }y=k:\quad
V=2\pi\int_c^d |y-k|\bigl(x_{\rm right}(y)-x_{\rm left}(y)\bigr)\,dy.
\end{align*}
\end{keyformula}

\subsubsection{Arc length and surface area}
\begin{keyformula}[title={Arc Length and Surface Area of Revolution}]
\begin{align*}
&L=\int_a^b\sqrt{1+(f'(x))^2}\,dx\quad\text{for }y=f(x),
\qquad
L=\int_c^d\sqrt{1+(g'(y))^2}\,dy\quad\text{for }x=g(y).\\
&S=2\pi\int_a^b |f(x)-k|\sqrt{1+(f'(x))^2}\,dx
\quad\text{about }y=k.\\
&S=2\pi\int_c^d |g(y)-k|\sqrt{1+(g'(y))^2}\,dy
\quad\text{about }x=k.
\end{align*}
\end{keyformula}

\subsubsection{Work and fluid force}
\begin{keyformula}[title={Work, Pumping, and Hydrostatic Force}]
\begin{align*}
&\textbf{Work:}\quad W=\int_a^b F(x)\,dx,
\qquad
\textbf{Hooke's law:}\quad F=kx,\\[-0.2em]
&\textbf{Pumping:}\quad dW=\rho g\,A(y)D(y)\,dy,
\qquad W=\int dW.\\[-0.2em]
&\textbf{Hydrostatic force:}\quad p(y)=\rho g h(y),
\qquad dF=p(y)w(y)\,dy,
\qquad F=\int dF.
\end{align*}
\end{keyformula}

\subsection{Ch3: Techniques of integration and numerical integration}

\subsubsection{Integration by parts}
\begin{keyformula}[title={Integration by Parts}]
\[
\int u\,dv=uv-\int v\,du,
\qquad
\int_a^b u\,dv=\bigl[uv\bigr]_a^b-\int_a^b v\,du.
\]
\end{keyformula}

\begin{examtip}
Choose \(u\) so that differentiating simplifies it.  A useful priority is:
logarithmic/inverse-trig, then algebraic, then trigonometric/exponential.
After one step, check that the remaining integral is simpler.
\end{examtip}

\subsubsection{Trigonometric integrals and substitutions}
\begin{keyformula}[title={Trigonometric Integral Decision Rules}]
\begin{align*}
&\int \sin^m x\cos^n x\,dx:
\quad n\text{ odd}\Rightarrow\text{save one }\cos x,\ u=\sin x;\quad
m\text{ odd}\Rightarrow\text{save one }\sin x,\ u=\cos x.\\
&\text{If both powers are even, use}
\quad \sin^2x=\frac{1-\cos2x}{2},\quad
\cos^2x=\frac{1+\cos2x}{2},\quad
\sin x\cos x=\frac12\sin2x.\\
&\int \tan^m x\sec^n x\,dx:
\quad n\text{ even}\Rightarrow\text{save }\sec^2x,\ u=\tan x;\quad
m\text{ odd}\Rightarrow\text{save }\sec x\tan x,\ u=\sec x.
\end{align*}
\end{keyformula}

\begin{keyformula}[title={Trigonometric Substitution Map}]
\[
\sqrt{a^2-x^2}\Rightarrow x=a\sin\theta,
\qquad
\sqrt{a^2+x^2}\Rightarrow x=a\tan\theta,
\qquad
\sqrt{x^2-a^2}\Rightarrow x=a\sec\theta.
\]
\end{keyformula}

\subsubsection{Partial fractions}
\begin{keyformula}[title={Partial Fraction Templates}]
\begin{align*}
&\text{Distinct linear factors:}\quad
\frac{P(x)}{\prod_i(a_i x+b_i)}=\sum_i\frac{A_i}{a_i x+b_i}.\\
&\text{Repeated linear factor }(ax+b)^r:\quad
\sum_{k=1}^r\frac{A_k}{(ax+b)^k}.\\
&\text{Irreducible quadratic factor }(ax^2+bx+c)^r:\quad
\sum_{k=1}^r\frac{A_kx+B_k}{(ax^2+bx+c)^k}.\\
&\text{Fast quadratic split:}\quad
Ax+B=\lambda(2ax+b)+\mu.
\end{align*}
\end{keyformula}

\begin{commonmistake}
If \(\deg P\ge \deg Q\), perform polynomial long division before partial
fractions.  An irreducible quadratic factor needs a linear numerator.
\end{commonmistake}

\subsubsection{Numerical integration}
Let \(\Delta x=(b-a)/n\), \(x_i=a+i\Delta x\).
\begin{keyformula}[title={Left, Right, Midpoint, Trapezoidal, Simpson Rules}]
\begin{align*}
&L_n=\sum_{i=1}^n f(x_{i-1})\Delta x,
\qquad
R_n=\sum_{i=1}^n f(x_i)\Delta x,
\qquad
M_n=\sum_{i=1}^n f\!\left(\frac{x_{i-1}+x_i}{2}\right)\Delta x.\\
&T_n=\frac{\Delta x}{2}\left(f(x_0)+2\sum_{i=1}^{n-1}f(x_i)+f(x_n)\right)=\frac{L_n+R_n}{2}.\\
&S_n=\frac{\Delta x}{3}\left(f(x_0)+4\sum_{\substack{1\le i\le n-1\\ i\text{ odd}}}f(x_i)+2\sum_{\substack{2\le i\le n-2\\ i\text{ even}}}f(x_i)+f(x_n)\right),
\quad n\text{ even}.
\end{align*}
\end{keyformula}

\begin{keyformula}[title={Numerical Error Bounds}]
\[
|E_T|\le \frac{K(b-a)^3}{12n^2},
\qquad
|E_M|\le \frac{K(b-a)^3}{24n^2}
\quad (|f''|\le K),
\]
\[
|E_S|\le \frac{K(b-a)^5}{180n^4}
\quad (|f^{(4)}|\le K,\ n\text{ even}).
\]
\end{keyformula}

\begin{examtip}
If \(f\) is increasing, \(L_n\le \int_a^b f\le R_n\); reverse for decreasing.
If \(f''>0\), trapezoidal overestimates and midpoint underestimates; reverse
for \(f''<0\).
\end{examtip}

\subsection{Ch4: Sequences and series}

\subsubsection{Sequences}
\begin{keyformula}[title={Sequence Convergence and Monotone Convergence}]
\[
a_n\to L\iff \forall\varepsilon>0\ \exists N:\ n\ge N\Rightarrow |a_n-L|<\varepsilon,
\qquad
\text{monotone and bounded}\Rightarrow\text{convergent}.
\]
\end{keyformula}

\subsubsection{Series basics}
\begin{keyformula}[title={Partial Sums, Remainders, and the $n$th-Term Test}]
\[
\sum_{n=1}^{\infty}a_n\text{ converges}\iff
S_n:=\sum_{k=1}^n a_k\to S\in\mathbb R,
\qquad
R_n:=S-S_n=\sum_{k=n+1}^{\infty}a_k.
\]
\[
\sum a_n\text{ convergent}\Rightarrow a_n\to0,
\qquad
a_n\not\to0\Rightarrow \sum a_n\text{ diverges}.
\]
\end{keyformula}

\subsubsection{Standard series and convergence tests}
\begin{keyformula}[title={Geometric and $p$-Series}]
\[
\sum_{n=0}^{\infty}ar^n=\frac{a}{1-r}\quad(|r|<1),
\qquad
\sum_{n=1}^{\infty}\frac{1}{n^p}\text{ converges}\Longleftrightarrow p>1.
\]
\end{keyformula}

\begin{keyformula}[title={Core Convergence Tests}]
Assume \(a_n\ge0\) where required.
\begin{align*}
&0\le a_n\le b_n,
\quad \sum b_n\text{ conv}\Rightarrow\sum a_n\text{ conv},
\quad \sum a_n\text{ div}\Rightarrow\sum b_n\text{ div}.\\
&a_n,b_n>0,
\quad \lim\frac{a_n}{b_n}=c\in(0,\infty)
\Rightarrow \sum a_n,\sum b_n\text{ have the same behaviour}.\\
&a_n=f(n),\ f>0,\ f\downarrow:
\quad \sum a_n\text{ conv}\Longleftrightarrow \int_1^\infty f(x)\,dx\text{ conv},
\quad
\int_{n+1}^\infty f\le R_n\le\int_n^\infty f.\\
&b_n\downarrow0:
\quad \sum (-1)^n b_n\text{ converges},
\qquad |R_n|\le b_{n+1}.\\
&\sum |a_n|\text{ conv}\Rightarrow\sum a_n\text{ conv absolutely};
\qquad
\sum a_n\text{ conv but }\sum|a_n|\text{ div}\Rightarrow\text{conditional}.\\
&L=\limsup\left|\frac{a_{n+1}}{a_n}\right|:
\quad L<1\Rightarrow\text{absolute convergence},
\quad L>1\Rightarrow\text{divergence}.\\
&L=\limsup\sqrt[n]{|a_n|}:
\quad L<1\Rightarrow\text{absolute convergence},
\quad L>1\Rightarrow\text{divergence}.
\end{align*}
\end{keyformula}

\begin{examtip}
Check \(a_n\to0\) first.  Then try geometric or \(p\)-series recognition,
ratio for factorials/exponentials, root for \(n\)-th powers, and comparison for
asymptotic matches.
\end{examtip}

\subsection{Ch5: Taylor / Maclaurin polynomials}

\subsubsection{Taylor polynomial and error}
\begin{keyformula}[title={Taylor Polynomial and Lagrange Remainder}]
\[
T_n(x)=\sum_{k=0}^{n}\frac{f^{(k)}(a)}{k!}(x-a)^k,
\qquad
f(x)=T_n(x)+R_n(x),
\]
\[
R_n(x)=\frac{f^{(n+1)}(\xi)}{(n+1)!}(x-a)^{n+1},
\qquad
|R_n(x)|\le \frac{M}{(n+1)!}|x-a|^{n+1}.
\]
\end{keyformula}

\subsubsection{Common Maclaurin series}
\begin{keyformula}[title={High-Frequency Maclaurin Series}]
\begin{align*}
&e^x=\sum_{n=0}^{\infty}\frac{x^n}{n!},
\qquad
\sin x=\sum_{n=0}^{\infty}(-1)^n\frac{x^{2n+1}}{(2n+1)!},
\qquad
\cos x=\sum_{n=0}^{\infty}(-1)^n\frac{x^{2n}}{(2n)!}.\\
&\frac{1}{1-x}=\sum_{n=0}^{\infty}x^n\quad(|x|<1),
\qquad
\ln(1+x)=\sum_{n=1}^{\infty}(-1)^{n+1}\frac{x^n}{n}\quad(-1<x\le1).\\
&\arctan x=\sum_{n=0}^{\infty}(-1)^n\frac{x^{2n+1}}{2n+1}
\quad(|x|\le1\text{ with endpoint care}).
\end{align*}
\end{keyformula}

\subsection{Ch6: Power series}

\subsubsection{Radius and interval of convergence}
\begin{keyformula}[title={Power Series Radius and Endpoint Workflow}]
\[
\sum_{n=0}^{\infty}c_n(x-a)^n
\quad\text{converges for }|x-a|<R\text{ and diverges for }|x-a|>R.
\]
\[
L(x)=\limsup_{n\to\infty}\left|\frac{c_{n+1}(x-a)^{n+1}}{c_n(x-a)^n}\right|<1
\quad\text{gives the open interval of convergence.}
\]
Endpoints \(x=a\pm R\) must be tested separately as numerical series.
\end{keyformula}

\subsubsection{Termwise operations}
\begin{keyformula}[title={Term-by-Term Differentiation and Integration}]
For \(|x-a|<R\),
\[
\frac{d}{dx}\sum_{n=0}^{\infty}c_n(x-a)^n
=\sum_{n=1}^{\infty}n c_n(x-a)^{n-1},
\]
\[
\int\sum_{n=0}^{\infty}c_n(x-a)^n\,dx
=C+\sum_{n=0}^{\infty}\frac{c_n}{n+1}(x-a)^{n+1}.
\]
The radius is unchanged; endpoint behaviour may change.
\end{keyformula}

\begin{commonmistake}
Do not use the radius test to decide endpoints.  Endpoint testing is a separate
ordinary-series problem.
\end{commonmistake}
